\documentclass[spanish,letterpaper]{article}
\usepackage[utf8]{inputenc}
\usepackage[T1]{fontenc}
\usepackage[spanish,es-nodecimaldot]{babel}
\usepackage{helvet}
\renewcommand{\familydefault}{\sfdefault}
\usepackage[margin=1.4cm,top=1.25cm,bottom=1.45cm]{geometry}
\usepackage{amsmath,amssymb,array,tabularx,colortbl,booktabs,tikz,fancyhdr}
\usetikzlibrary{arrows.meta,calc,angles,quotes,babel}
\usepackage{microtype}
\usepackage{xcolor}

\definecolor{iiiepeNavy}{HTML}{163B5C}
\definecolor{iiiepeBlue}{HTML}{1F6F9F}
\definecolor{iiiepeAqua}{HTML}{20B7A8}
\definecolor{iiiepePale}{HTML}{EEF5F8}
\definecolor{iiiepeInk}{HTML}{1E2A32}
\definecolor{okGreen}{HTML}{1B6B3A}
\definecolor{errRed}{HTML}{A61C1C}

\setlength{\parindent}{0pt}
\setlength{\parskip}{4pt}
\renewcommand{\arraystretch}{1.28}
\arrayrulecolor{iiiepeNavy!65}

\pagestyle{fancy}
\fancyhf{}
\renewcommand{\headrulewidth}{0pt}
\fancyfoot[L]{\scriptsize IIIEPE · MAEM · PFM02 · SD26}
\fancyfoot[C]{\scriptsize\thepage}
\fancyfoot[R]{\scriptsize Parcial 1 · Sospechoso, muy sospechoso}

\begin{document}
\color{iiiepeInk}

\noindent
\begin{tabularx}{\textwidth}{@{} >{\raggedright\arraybackslash}p{0.44\textwidth} >{\centering\arraybackslash}p{0.31\textwidth} >{\raggedleft\arraybackslash}X @{}}
\textbf{\footnotesize Maestría en Aprendizaje y Enseñanza de las Matemáticas} &
\textbf{\footnotesize UA: Temas Selectos de Matemáticas I} &
\textbf{\footnotesize Bloque II. Enseñar desde el desafío} \\
\footnotesize \textbf{Parcial 1. Sospechoso, muy sospechoso} &
\footnotesize Sesión 10. Universo 3-60 &
\footnotesize Evaluación: \rule{1.8cm}{0.4pt}
\end{tabularx}

\vspace{4pt}
\noindent
\begin{tabularx}{\textwidth}{@{} >{\raggedright\arraybackslash}X >{\raggedleft\arraybackslash}p{0.34\textwidth} @{}}
\textbf{Nombre:} \underline{\hspace{0.2cm}Mtro. Martín Jonathan de la Cruz Muñoz (Matrícula 00874)\hspace{0.2cm}} &
\textbf{Fecha:} \underline{\hspace{0.2cm}24 de octubre de 2026\hspace{0.2cm}}
\end{tabularx}

\vspace{4pt}
{\small\textbf{Instrucciones:} Lee cuidadosamente cada afirmación. Indica si es cierta o falsa y justifica tu respuesta mediante un argumento matemático, un ejemplo o un contraejemplo. No basta con marcar una opción.}

\vspace{4pt}

\noindent\begin{tabularx}{\textwidth}{|>{\raggedright\arraybackslash}p{0.25\textwidth}|>{\centering\arraybackslash}p{0.06\textwidth}|>{\centering\arraybackslash}p{0.06\textwidth}|>{\raggedright\arraybackslash}X|}
\hline
\rowcolor{iiiepePale}
\multicolumn{1}{|c|}{\textbf{Afirmación}} & \textbf{Cierta} & \textbf{Falsa} & \multicolumn{1}{c|}{\textbf{Justificación (argumento, ejemplo o contraejemplo)}} \\\hline

1. Entre dos números decimales distintos siempre es posible encontrar otro número decimal. &
\textbf{\large\color{okGreen}X} &
&
\textbf{Verdadera (propiedad de densidad).} Si $a < b$ son dos decimales distintos, su punto medio $m = \frac{a+b}{2}$ cumple $a < m < b$ y sigue siendo decimal exacto (como $2$ divide a $10$, dividir entre $2$ a lo sumo añade una cifra decimal más). Además, se pueden generar infinitos entre $a$ y $b$.\par\smallskip
\textit{Ejemplo:} Entre $1.4$ y $1.5$ está $m = \frac{1.4+1.5}{2} = 1.45$ ($1.4 < 1.45 < 1.5$); entre $1.45$ y $1.46$ está $1.455$. \\\hline

2. Si a dos fracciones positivas les sumamos la misma fracción positiva, la diferencia entre sus valores aumenta. &
&
\textbf{\large\color{errRed}X} &
\textbf{Falsa (invariancia por traslación aditiva).} Sean $p, q, r > 0$ con $p > q$. Al sumar $r$ a ambas fracciones, la nueva diferencia es:
\[
(p + r) - (q + r) = p - q + r - r = p - q.
\]
La distancia entre los valores permanece \textbf{constante}; no aumenta ni disminuye.\par
\textit{Contraejemplo:} Para $p = \frac{3}{4}$ y $q = \frac{1}{4}$, su diferencia es $\frac{3}{4}-\frac{1}{4}=\frac{1}{2}$. Si sumamos $r = \frac{1}{2}$, resultan $\frac{5}{4}$ y $\frac{3}{4}$, cuya diferencia sigue siendo $\frac{5}{4}-\frac{3}{4}=\frac{1}{2}$. \\\hline

3. Si sumamos el mismo número a ambos lados de una ecuación, las soluciones de la ecuación no cambian. &
\textbf{\large\color{okGreen}X} &
&
\textbf{Verdadera (propiedad uniforme aditiva de la igualdad).} Para todo $c \in \mathbb{R}$, las ecuaciones $f(x) = g(x)$ y $f(x) + c = g(x) + c$ son \textbf{equivalentes} ($f(x)=g(x) \iff f(x)+c=g(x)+c$), pues la operación es reversible sumando el inverso aditivo $-c$. Por tanto, el conjunto solución no se altera.\par\smallskip
\textit{Ejemplo:} En $2x - 5 = 7$, la única solución es $x = 6$. Al sumar $+5$ a ambos lados se obtiene $2x = 12$, cuya única solución continúa siendo $x = 6$. \\\hline

4. Si un punto se encuentra sobre la mediatriz de un segmento, entonces necesariamente está dentro del segmento. &
&
\textbf{\large\color{errRed}X} &
\textbf{Falsa.} La mediatriz de $\overline{AB}$ es la \textbf{recta infinita} perpendicular a $\overline{AB}$ que pasa por su punto medio $M$, caracterizada por $\{P : PA = PB\}$. De todos sus infinitos puntos, únicamente $M$ pertenece al segmento $\overline{AB}$; cualquier otro punto $P \neq M$ está \textbf{fuera} del segmento.\par\smallskip
\textit{Contraejemplo:} Para $A=(0,0)$ y $B=(4,0)$, la mediatriz es $x=2$. El punto $P=(2,3)$ pertenece a ella ($PA=PB=\sqrt{13}$), pero no está contenido en el segmento $\overline{AB}$ (donde $y=0$). \\\hline

5. Si un punto se encuentra sobre la bisectriz interior de un ángulo, entonces está a la misma distancia de los dos lados del ángulo. &
\textbf{\large\color{okGreen}X} &
&
\textbf{Verdadera (caracterización de la bisectriz).} Sea $P$ un punto en la bisectriz de $\angle AOB$, y sean $\overline{PH}\perp \overrightarrow{OA}$ y $\overline{PK}\perp \overrightarrow{OB}$ las distancias perpendiculares a cada lado. Los triángulos rectángulos $\triangle OHP$ y $\triangle OKP$ comparten la hipotenusa $OP$, tienen $\angle HOP = \angle KOP$ y $\angle OHP = \angle OKP = 90^\circ$; por criterio AAL, $\triangle OHP \cong \triangle OKP \implies PH = PK$.\par\smallskip
\textit{Ejemplo:} En el primer cuadrante, $P=(4,4)$ sobre la bisectriz $y=x$ dista exactamente $4\text{ u}$ tanto del eje $x$ como del eje $y$. \\\hline
\end{tabularx}

\vspace{8pt}
\noindent
\begin{minipage}[t]{0.48\textwidth}
  \centering
  \begin{tikzpicture}[scale=0.82]
    \draw[fill=iiiepePale!60,draw=iiiepeBlue!50,rounded corners=3pt] (-0.6,-0.75) rectangle (5.6,3.75);
    \node[font=\scriptsize\bfseries,iiiepeNavy] at (2.5,3.45) {Contraejemplo Afirmación 4 ($P \notin \overline{AB}$)};
    % Segmento AB
    \draw[very thick,iiiepeNavy] (0.5,0.4) -- (4.5,0.4);
    \draw[iiiepeAqua!80!black,thick,dashed] (2.5,-0.45) -- (2.5,3.05);
    \draw (2.5,0.4) rectangle +(0.25,0.25);
    % Segmentos PA y PB
    \draw[errRed,thick,dotted] (2.5,2.6) -- (0.5,0.4);
    \draw[errRed,thick,dotted] (2.5,2.6) -- (4.5,0.4);
    \fill[iiiepeNavy] (0.5,0.4) circle (2.2pt) node[below,font=\scriptsize] {$A(0,0)$};
    \fill[iiiepeNavy] (4.5,0.4) circle (2.2pt) node[below,font=\scriptsize] {$B(4,0)$};
    \fill[iiiepeAqua!70!black] (2.5,0.4) circle (2.2pt) node[below right,font=\scriptsize] {$M(2,0)$};
    \fill[errRed] (2.5,2.6) circle (2.6pt) node[above right,font=\scriptsize\bfseries] {$P(2,3) \notin \overline{AB}$};
    \node[font=\tiny,errRed,rotate=48] at (1.3,1.65) {$\sqrt{13}$};
    \node[font=\tiny,errRed,rotate=-48] at (3.7,1.65) {$\sqrt{13}$};
    \node[font=\tiny,iiiepeAqua!70!black,fill=iiiepePale!60,inner sep=1pt] at (3.35,0.85) {Mediatriz $x=2$};
  \end{tikzpicture}
\end{minipage}\hfill
\begin{minipage}[t]{0.48\textwidth}
  \centering
  \begin{tikzpicture}[scale=0.82]
    \draw[fill=iiiepePale!60,draw=iiiepeBlue!50,rounded corners=3pt] (-0.5,-0.75) rectangle (5.7,3.75);
    \node[font=\scriptsize\bfseries,iiiepeNavy] at (2.6,3.45) {Demostración Afirmación 5 ($PH = PK$)};
    \coordinate (O) at (0.3,0.2);
    \coordinate (A) at (4.9,0.2);
    \coordinate (B) at (3.1,2.82);
    \coordinate (P) at (3.7,1.54);
    \coordinate (H) at (3.7,0.2);
    \coordinate (K) at (2.78,2.52);
    \draw[-{Latex[length=2mm]},very thick,iiiepeNavy] (O) -- (A) node[below,font=\scriptsize] {$A$};
    \draw[-{Latex[length=2mm]},very thick,iiiepeNavy] (O) -- (B) node[below right,font=\scriptsize] {$B$};
    \draw[-{Latex[length=2mm]},thick,iiiepeAqua!80!black,dashed] (O) -- (4.55,1.88);
    \node[font=\tiny\bfseries,iiiepeAqua!70!black] at (4.85,1.72) {Bisectriz};
    \draw[okGreen,very thick] (P) -- (H);
    \draw[okGreen,very thick] (P) -- (K);
    \draw (H) rectangle +(-0.22,0.22);
    \fill[iiiepeNavy] (O) circle (2pt) node[below left,font=\scriptsize] {$O$};
    \fill[okGreen] (P) circle (2.5pt) node[above right,font=\scriptsize\bfseries] {$P$};
    \fill[iiiepeNavy] (H) circle (2pt) node[below,font=\scriptsize] {$H$};
    \fill[iiiepeNavy] (K) circle (2pt) node[above left,font=\scriptsize] {$K$};
    \node[font=\tiny\bfseries,okGreen,right] at (3.73,0.85) {$PH=PK$};
    \node[font=\tiny,iiiepeInk] at (2.8,-0.38) {$\triangle OHP \cong \triangle OKP$ (criterio AAL)};
  \end{tikzpicture}
\end{minipage}

\end{document}
